% GENERATED FILE. Edit canonical lessons, then run npm run generate:books.
% canonical-source-hash: fnv1a64:0240ea4de32fadfc
% canonical-lessons: ML-C515-ketupatu, ML-C515-nastapariharam, ML-C515-buddhimuttu, ML-C515-kurunnuka, ML-C515-tappuka

\chapter{Damage and Hardship}
\label{ch:ml-damage-and-hardship}
\hlchaptermodality{515}

\begin{chapteropening}
\textbf{By the end of this chapter:} \emph{I can say breakage, compensation, a hardship, to get tangled and to grope for.}

Complete the last lesson of chapter 515: I can say breakage, compensation, a hardship, to get tangled and to grope for.
\end{chapteropening}

\section[\texorpdfstring{\ml{കേടുപാട്} (kēṭupāṭŭ)}{കേടുപാട് (kēṭupāṭŭ)}]{\ml{കേടുപാട്} (kēṭupāṭŭ) — breakage; damage}

\label{lesson:ML-C515-ketupatu}



\begin{quote}
Before the new one: say the Malayalam for to be disconnected, then the Malayalam for to stop working.
\end{quote}

\subsection*{You'll want to know: \ml{കേടുപാട്}}
\textbf{\ml{കേടുപാട്}} — \emph{kēṭupāṭŭ} — \textquotedblleft{}breakage; damage\textquotedblright{}.

\begin{quote}
\ml{വെള്ളപ്പൊക്കത്തിൽ} \ml{വീടിന്} \ml{കേടുപാട്} \ml{പറ്റി}.
\end{quote}

\begin{quote}
\emph{veḷḷappokkattil vīṭinŭ kēṭupāṭŭ paṟṟi.}
\end{quote}

\begin{quote}
The house was damaged in the flood.
\end{quote}

\ml{കേട്}, damage, plus \ml{പാട്}, a mark. Usually with \ml{പറ്റുക}: \ml{കേടുപാട്} \ml{പറ്റി}.

\begin{cousinweb}
\textbf{\ml{കേടുപാട്}} (\emph{kēṭupāṭŭ}) is a native Dravidian compound of \ml{കേട്}, damage or loss (DEDR 1942), and \ml{പാട്}, a mark.
\end{cousinweb}

\begin{grammarlens}[title={What you've built}]
One more word for getting around.
\end{grammarlens}

\subsection*{Your turn}
\begin{itemize}
\raggedright
  \item \emph{Say it:} \emph{kēṭupāṭŭ}
  \item \emph{Say it:} \emph{kēṭupāṭŭ}, once more
  \item \emph{Say it:} \emph{veḷḷappokkattil vīṭinŭ kēṭupāṭŭ paṟṟi.} — The house was damaged in the flood.
  \item \emph{Recall:} say the Malayalam for to be disconnected, then the Malayalam for to stop working, then say \emph{kēṭupāṭŭ} again
\end{itemize}

\begin{tcolorbox}[breakable,colback=teal!4,colframe=teal!35!black,title={Before you move on}]
What does \ml{കേടുപാട്} mean? (\textquotedblleft{}Breakage; damage\textquotedblright{}.) Say it once more.
\end{tcolorbox}
\section[\texorpdfstring{\ml{നഷ്ടപരിഹാരം} (naṣṭaparihāraṁ)}{നഷ്ടപരിഹാരം (naṣṭaparihāraṁ)}]{\ml{നഷ്ടപരിഹാരം} (naṣṭaparihāraṁ) — compensation}

\label{lesson:ML-C515-nastapariharam}



\begin{quote}
Before the new one: say the Malayalam for to stop working, then the Malayalam for breakage.
\end{quote}

\subsection*{You'll want to know: \ml{നഷ്ടപരിഹാരം}}
\textbf{\ml{നഷ്ടപരിഹാരം}} — \emph{naṣṭaparihāraṁ} — \textquotedblleft{}compensation\textquotedblright{}.

\begin{quote}
\ml{വിമാനക്കമ്പനി} \ml{യാത്രക്കാർക്ക്} \ml{നഷ്ടപരിഹാരം} \ml{നൽകി}.
\end{quote}

\begin{quote}
\emph{vimānakkampani yātrakkārkkŭ naṣṭaparihāraṁ nalki.}
\end{quote}

\begin{quote}
The airline paid the passengers compensation.
\end{quote}

\ml{നഷ്ടം}, a loss, plus \ml{പരിഹാരം}, a remedy: a remedy for a loss.

\begin{cousinweb}
\textbf{\ml{നഷ്ടപരിഹാരം}} (\emph{naṣṭaparihāraṁ}) is a Malayalam compound built from two Sanskrit words, \emph{naṣṭa}, \textquotedblleft{}lost\textquotedblright{} (CDIAL 7027), and \emph{parihāra}, \textquotedblleft{}avoiding\textquotedblright{} (CDIAL 7901), which Monier-Williams also gives as removing and remedying.
\end{cousinweb}

\begin{grammarlens}[title={What you've built}]
One more word for getting around.
\end{grammarlens}

\subsection*{Your turn}
\begin{itemize}
\raggedright
  \item \emph{Say it:} \emph{naṣṭaparihāraṁ}
  \item \emph{Say it:} \emph{naṣṭaparihāraṁ}, once more
  \item \emph{Say it:} \emph{vimānakkampani yātrakkārkkŭ naṣṭaparihāraṁ nalki.} — The airline paid the passengers compensation.
  \item \emph{Recall:} say the Malayalam for to stop working, then the Malayalam for breakage, then say \emph{naṣṭaparihāraṁ} again
\end{itemize}

\begin{tcolorbox}[breakable,colback=teal!4,colframe=teal!35!black,title={Before you move on}]
What does \ml{നഷ്ടപരിഹാരം} mean? (\textquotedblleft{}Compensation\textquotedblright{}.) Say it once more.
\end{tcolorbox}
\section[\texorpdfstring{\ml{ബുദ്ധിമുട്ട്} (buddhimuṭṭŭ)}{ബുദ്ധിമുട്ട് (buddhimuṭṭŭ)}]{\ml{ബുദ്ധിമുട്ട്} (buddhimuṭṭŭ) — a hardship; difficulty}

\label{lesson:ML-C515-buddhimuttu}



\begin{quote}
Before the new one: say the Malayalam for breakage, then the Malayalam for compensation.
\end{quote}

\subsection*{You'll want to know: \ml{ബുദ്ധിമുട്ട്}}
\textbf{\ml{ബുദ്ധിമുട്ട്}} — \emph{buddhimuṭṭŭ} — \textquotedblleft{}a hardship; difficulty\textquotedblright{}.

\begin{quote}
\ml{വെള്ളം} \ml{ഇല്ലാത്തതുകൊണ്ട്} \ml{ഞങ്ങൾക്ക്} \ml{വലിയ} \ml{ബുദ്ധിമുട്ടുണ്ടായി}.
\end{quote}

\begin{quote}
\emph{veḷḷaṁ illāttatukoṇṭŭ ñaṅṅaḷkkŭ valiya buddhimuṭṭuṇṭāyi.}
\end{quote}

\begin{quote}
We had great difficulty because there was no water.
\end{quote}

The one who has the difficulty takes -\ml{ക്ക്}: \ml{എനിക്ക്} \ml{ബുദ്ധിമുട്ടുണ്ട്}. \ml{ബുദ്ധിമുട്ടുക} is to struggle.

\begin{cousinweb}
\textbf{\ml{ബുദ്ധിമുട്ട്}} (\emph{buddhimuṭṭŭ}) is a Malayalam compound built from \ml{ബുദ്ധി}, the mind, from Sanskrit \emph{buddhi} (Monier-Williams), and the native \ml{മുട്ട്}, an impediment or stoppage (DEDR 4933).
\end{cousinweb}

\begin{grammarlens}[title={What you've built}]
One more word for getting around.
\end{grammarlens}

\subsection*{Your turn}
\begin{itemize}
\raggedright
  \item \emph{Say it:} \emph{buddhimuṭṭŭ}
  \item \emph{Say it:} \emph{buddhimuṭṭŭ}, once more
  \item \emph{Say it:} \emph{veḷḷaṁ illāttatukoṇṭŭ ñaṅṅaḷkkŭ valiya buddhimuṭṭuṇṭāyi.} — We had great difficulty because there was no water.
  \item \emph{Recall:} say the Malayalam for breakage, then the Malayalam for compensation, then say \emph{buddhimuṭṭŭ} again
\end{itemize}

\begin{tcolorbox}[breakable,colback=teal!4,colframe=teal!35!black,title={Before you move on}]
What does \ml{ബുദ്ധിമുട്ട്} mean? (\textquotedblleft{}A hardship; difficulty\textquotedblright{}.) Say it once more.
\end{tcolorbox}
\section[\texorpdfstring{\ml{കുരുങ്ങുക} (kuruṅṅuka)}{കുരുങ്ങുക (kuruṅṅuka)}]{\ml{കുരുങ്ങുക} (kuruṅṅuka) — to get tangled; to get caught}

\label{lesson:ML-C515-kurunnuka}



\begin{quote}
Before the new one: say the Malayalam for compensation, then the Malayalam for a hardship.
\end{quote}

\subsection*{You'll want to know: \ml{കുരുങ്ങുക}}
\textbf{\ml{കുരുങ്ങുക}} — \emph{kuruṅṅuka} — \textquotedblleft{}to get tangled; to get caught\textquotedblright{}.

\begin{quote}
\ml{പട്ടത്തിന്റെ} \ml{നൂല്} \ml{മരത്തിൽ} \ml{കുരുങ്ങി}.
\end{quote}

\begin{quote}
\emph{paṭṭattinṟe nūlŭ marattil kuruṅṅi.}
\end{quote}

\begin{quote}
The kite string got tangled in the tree.
\end{quote}

Past: \ml{കുരുങ്ങി}, \emph{kuruṅṅi}.

What it gets caught in takes -\ml{ഇൽ}: \ml{വലയിൽ} \ml{കുരുങ്ങി}, caught in the net.

\begin{cousinweb}
\textbf{\ml{കുരുങ്ങുക}} (\emph{kuruṅṅuka}) is a Malayalam verb related to \ml{കുരുക്ക്}, a knot or noose; its origin is not settled.
\end{cousinweb}

\begin{grammarlens}[title={What you've built}]
One more word for getting around.
\end{grammarlens}

\subsection*{Your turn}
\begin{itemize}
\raggedright
  \item \emph{Say it:} \emph{kuruṅṅuka}
  \item \emph{Say it:} \emph{kuruṅṅuka}, once more
  \item \emph{Say it:} \emph{paṭṭattinṟe nūlŭ marattil kuruṅṅi.} — The kite string got tangled in the tree.
  \item \emph{Recall:} say the Malayalam for compensation, then the Malayalam for a hardship, then say \emph{kuruṅṅuka} again
\end{itemize}

\begin{tcolorbox}[breakable,colback=teal!4,colframe=teal!35!black,title={Before you move on}]
What does \ml{കുരുങ്ങുക} mean? (\textquotedblleft{}To get tangled; to get caught\textquotedblright{}.) Say it once more.
\end{tcolorbox}
\section[\texorpdfstring{\ml{തപ്പുക} (tappuka)}{തപ്പുക (tappuka)}]{\ml{തപ്പുക} (tappuka) — to grope for; to feel about}

\label{lesson:ML-C515-tappuka}



\begin{quote}
Before the new one: say the Malayalam for a hardship, then the Malayalam for to get tangled.
\end{quote}

\subsection*{You'll want to know: \ml{തപ്പുക}}
\textbf{\ml{തപ്പുക}} — \emph{tappuka} — \textquotedblleft{}to grope for; to feel about\textquotedblright{}.

\begin{quote}
\ml{ഇരുട്ടിൽ} \ml{ഞാൻ} \ml{ടോർച്ചിനായി} \ml{തപ്പി}.
\end{quote}

\begin{quote}
\emph{iruṭṭil ñān ṭōrccināyi tappi.}
\end{quote}

\begin{quote}
In the dark I groped for the torch.
\end{quote}

Past: \ml{തപ്പി}, \emph{tappi}.

What you grope for takes -\ml{ആയി}: \ml{ടോർച്ചിനായി} \ml{തപ്പി}.

\begin{cousinweb}
\textbf{\ml{തപ്പുക}} (\emph{tappuka}) is native Dravidian, \textquotedblleft{}to grope, to feel about\textquotedblright{} (DEDR 3072).
\end{cousinweb}

\begin{grammarlens}[title={What you've built}]
One more word for getting around.
\end{grammarlens}

\subsection*{Your turn}
\begin{itemize}
\raggedright
  \item \emph{Say it:} \emph{tappuka}
  \item \emph{Say it:} \emph{tappuka}, once more
  \item \emph{Say it:} \emph{iruṭṭil ñān ṭōrccināyi tappi.} — In the dark I groped for the torch.
  \item \emph{Recall:} say the Malayalam for a hardship, then the Malayalam for to get tangled, then say \emph{tappuka} again
\end{itemize}

\begin{tcolorbox}[breakable,colback=teal!4,colframe=teal!35!black,title={Before you move on}]
What does \ml{തപ്പുക} mean? (\textquotedblleft{}To grope for; to feel about\textquotedblright{}.) Say it once more.
\end{tcolorbox}
